
\subsection{Likelihood Estimation}
\label{sec:likelihood-estimator-variance}

The analytic integrals in Equation~\eqref{eq:hierarchical_likelihood} are not tractable, and so we estimate the integrals using Monte Carlo estimation \citep{Tiwari:2017ndi, Farr:2019rap, Essick:2022ojx, Talbot:2023pex}. 
For example, the estimator for the likelihood $\hat{\mathcal{L}}(d_i|\PEhyperparam)$ is
\begin{equation}
    \hat{\mathcal{L}}(d_i|\PEhyperparam) \propto \frac{1}{\Npe} \sum_{j=1}^{\Npe} \frac{\pi(\PEparameter_{ij}|\PEhyperparam)}{p(\PEparameter_{ij})},
    \label{eq:single_event_estimators}
\end{equation}
where $\{\PEparameter_{ij}\}_{j=1}^{\Npe}$ are a collection of $\Npe$ samples from the posterior on the \ac{GW} parameters of the $i{\rm th}$ event $d_i$. 
We divide out by the parameter estimation prior $p(\PEparameter)$, and so
the Monte Carlo estimator in Equation~\eqref{eq:single_event_estimators} converges to the $i{\rm th}$ integral inside the product of Equation~\eqref{eq:hierarchical_likelihood} in the limit of $\Npe \to \infty$.

Similarly, the estimator for the selection efficiency $\hat{\xi}$ is \citep{Tiwari:2017ndi, Farr:2019rap, Essick:2022ojx}
\begin{equation}
    \hat{\xi}(\PEhyperparam) \propto \frac{1}{\Ndraw} \sum_{j=1}^{\Nfound} \frac{\pi(\PEparameter_{j}|\PEhyperparam)}{\pi(\PEparameter_{j}|\PEhyperparam_{\rm draw})},
    \label{eq:selection_estimator}
\end{equation}
where $\Ndraw$ events with parameters $\PEparameter_i$ are injected into representative noise from the detectors. 
The search pipelines then search these synthetically generated data and recover some subset $\Nfound$ of the events with the detection statistic exceeding some threshold.
In the limit of $\Ndraw \to \infty$, this approaches the true integral in Equation~\eqref{eq:selection_efficiency}.

Because we only have a finite number of samples from each event and finite $\Ndraw$, we must be careful to account for the intrinsic variance in the estimation of the likelihood. 
To be sure our Monte Carlo estimators for the likelihood are trustworthy, we study the impact of Monte Carlo uncertainty in every inference we perform. 
Specifically, we compute the variance in the log-likelihood estimator, which varies across parameter space due to our resampling techniques in Equation~\eqref{eq:single_event_estimators} and Equation~\eqref{eq:selection_estimator}.
Propagating the uncertainty in the log-likelihood along independent degrees of freedom, the variance in the log-likelihood estimator $\sigma^2_{\ln \hat{\mathcal{L}}}$ in combining Equation~\eqref{eq:hierarchical_likelihood}, Equation~\eqref{eq:selection_efficiency}, and Equation~\eqref{eq:single_event_estimators} can be estimated as~(e.g., \citealt{Essick:2022ojx})
\begin{equation}
	\sigma^2_{\ln \hat{\mathcal{L}}}(\PEhyperparam) = \sum_{i=1}^{N_{\rm det}}\frac{\sigma^2_{\hat{\mathcal{L}}_i}(\PEhyperparam)}{\hat{\mathcal{L}}_i(\PEhyperparam)^2} + N(\Lambda)^2 \sigma^2_{\xi}(\PEhyperparam) ,
	\label{eq:variance_likelihood}
\end{equation}
where
\begin{equation}
	\sigma^2_{\hat{\mathcal{L}}_i}(\PEhyperparam) = \frac{1}{N_{\rm PE}}\left[\frac{1}{N_{\rm PE} -1 }\sum_{j=1}^{N_{\rm PE}}\left(\frac{\pi(\PEparameter_{ij}|\PEhyperparam)}{p(\PEparameter_{ij})}\right)^2 - \hat{\mathcal{L}}_i(\PEhyperparam)^2 \right]
\end{equation}
is the Monte Carlo variance in the single event Monte Carlo integrals of Equation~\eqref{eq:single_event_estimators} and 
\begin{equation}
	\sigma^2_{\xi}(\PEhyperparam) = \frac{1}{N_{\rm draw}}\left[\frac{1}{N_{\rm draw} -1}\sum_{j=1}^{N_{\rm found}}\left(\frac{\pi(\PEparameter_{j}|\PEhyperparam)}{p(\PEparameter_{j}|\PEhyperparam_{\rm draw})}\right)^2 - \hat{\xi}(\PEhyperparam)^2\right] 
\end{equation}
is the variance in the detection efficiency Monte Carlo integral of Equation~\eqref{eq:selection_estimator}.
When the rate-marginalized likelihood of Equation~\eqref{eq:rate_marginalized_hierarchical_likelihood} is used, $\sigma^2_{\ln \hat{\mathcal{L}}}$ takes a slightly different form
\begin{equation}
	\sigma^2_{\ln \hat{\mathcal{L}}}(\PEhyperparam) = \sum_{i=1}^{N_{\rm det}}\frac{\sigma^2_{\hat{\mathcal{L}}_i}(\PEhyperparam)}{\hat{\mathcal{L}}_i(\PEhyperparam)^2} + N_{\rm det}^2 \frac{\sigma^2_{\xi}(\PEhyperparam)}{\hat{\xi}(\PEhyperparam)^2}.
	\label{eq:variance_rate_marginalized_likelihood}
\end{equation}

It has been shown that \ac{GW} population inference can be biased when the variance in log-likelihood estimator exceeds $1$.
Consequently, we adopt a threshold on $\sigma^2_{\ln \hat{\mathcal{L}}}$ of $1$ to manage the bias of the posterior. Above this threshold the likelihood estimate may not be converged, and thus we ignore posterior samples with variances above this chosen threshold. Equation~\eqref{eq:variance_likelihood} describes the pointwise variance in the estimation of the log-likelihood. However, for accurate sampling of the posterior we only require that the difference of log-likelihoods to be small. A small pointwise log-likelihood variance is \textit{sufficient} for the variance of the difference of log-likelihoods to be small but not \textit{necessary}, rendering our threshold conservative~\citep{Farr:2019rap, Essick:2022ojx}. Indeed, for some models, a large region of hyperparameter space is removed by this threshold, limiting the range of potential populations that can be explored; for an example, see Appendix~\ref{appendix:model_comparison_effective_spins}.
Improvements to likelihood estimation are an active area of ongoing research~\citep{Wysocki:2018mpo,Doctor:2019ruh,Delfavero:2021qsc,Golomb:2021tll,Mould:2023eca,Hussain:2024qzl,Mancarella:2025uat}.

\subsection{Sampling Techniques}

In each model, we draw samples from the posterior to study the population distributions consistent with the data and the population model.
We draw samples using a variety of stochastic sampling algorithms, where the exact approach depends on the model. 
For most \parametrics, we use the nested sampler \DYNESTY\citep{Speagle:2019ivv} wrapper in \BILBY\citep{Ashton:2018jfp}, and the \GWPOPULATION implementation of the hierarchical likelihood \citep{Talbot:2024yqw}.

However, our \nonparametrics tend to have a large number of hyperparameters and so have a high dimensional posterior.
Nested sampling struggles with high dimensional distributions, so we use the \ac{HMC} adaptive \ac{NUTS} implementation in \NUMPYRO\citep{Phan:2019elc, Bingham:2019}. 
The \NUMPYRO adaptive \ac{NUTS} requires an autodifferentiable implementation of the likelihood, which we write in \JAX\citep{Bradbury:2018}.
This gradient information allows the \ac{NUTS} algorithm to efficiently explore high dimensional posteriors. 

